\documentclass[a4paper,12pt]{ctexbook}
\usepackage[margin=1cm]{geometry}
\usepackage{graphicx}
\usepackage{subfigure}
\usepackage{float}
\usepackage[colorlinks,linkcolor=black]{hyperref}%colorlinks启用链接颜色，linkcolor指定对应的颜色
\usepackage{listings}
%\usepackage[cache=false]{minted}  % 代码高亮X
\usepackage{fontspec}
\usepackage{tikz,xcolor,mwe}

\definecolor{cvgreen}{HTML}{92D14F}
\definecolor{cvgray}{HTML}{D8E4BE}
\definecolor{cvtext}{HTML}{92909B}
\usetikzlibrary{shadows}

\pagestyle{empty}
\CTEXsetup[format={\large\bfseries}]{section}%可以让section的标题左对齐。
\CTEXsetup[format={\Large\bfseries}]{chapter}%可以让chapter的标题左对齐。
%\CTEXsetup[format+={\flushleft}]{section}%让section的标题居左
%\renewcommand{\thesection}{\chinese{section}}%将“1.1”改为汉字“一”，但是subsection就会变成  六.1 ，比较难看，还是不用比较好。

\begin{document}


\begin{tikzpicture}[remember picture,overlay]
\fill[cvgreen] (current page.north west) rectangle ([xshift=6cm]current page.south west);% green bar
\fill[cvgray] ([yshift=-11cm]current page.north west) rectangle ([yshift=-17cm]current page.north east); % gray bar

\node[cvtext,right] at ([xshift=3.5cm,yshift=-13cm]current page.north west) {\Huge My ML way};
\node[cvtext,above left] at ([xshift=-1cm,yshift=-16.5cm]current page.north east) {\Large\bfseries \today};
% cover photo
\node[inner sep=0pt,below right] (image) at ([xshift=17cm,yshift=-1cm]current page.north west) {\includegraphics[width=3cm]{hainulogo.png}};
% name and address
\node[fill=white,align=center,text width=6.4cm,inner sep=0.8cm,below] (name) at (image.south) {};
\node[text width=15cm,inner sep=0.3cm,below right] at (name.south west){\Large Author:Flynn\\E-Mail:zofon@qq.com};
% attachments
\node[white,text width=5cm,inner sep=0.6cm,above right] at ([yshift=1cm]current page.south west)
{\large\obeylines\textbf{This is\\Open Source}};
\end{tikzpicture}

\tableofcontents

\chapter{相关的概念}
在学习机器学习之前首先要了解一些相关的概念。
\begin{description}
  \item[AI] 人工智能是最早出现的，也是最大、最外侧的同心圆；
  \item[ML] 机器学习
  \item[DL] 深度学习，深度学习造成了前所未有的巨大的影响， 是当今人工智能大爆炸的核心驱动。
\end{description}
五十年代，人工智能曾一度被极为看好。之后，人工智能的一些较小的子集发展了起来。先是机器学习，然后是深度学习。深度学习又是机器学习的子集。。

早在1956年夏天那次会议，人工智能的先驱们就梦想着用当时刚刚出现的计算机来构造复杂的、拥有与人类智慧同样本质特性的机器。这就是我们现在所说的“强人工智能”（General AI）。这个无所不能的机器，它有着我们所有的感知（甚至比人更多），我们所有的理性，可以像我们一样思考。

人们在电影里也总是看到这样的机器：友好的，像星球大战中的C-3PO；邪恶的，如终结者。强人工智能现在还只存在于电影和科幻小说中，原因不难理解，我们还没法实现它们，至少目前还不行。

我们目前能实现的，一般被称为“弱人工智能”（Narrow AI）。弱人工智能是能够与人一样，甚至比人更好地执行特定任务的技术。例如，Pinterest上的图像分类；或者Facebook的人脸识别。

这些是弱人工智能在实践中的例子。这些技术实现的是人类智能的一些具体的局部。但它们是如何实现的？这种智能是从何而来？这就带我们来到同心圆的里面一层，机器学习。

\chapter{如何研究学习一个机器学习算法}
中文：\url{http://blog.jobbole.com/80658/}
英文：\url{http://machinelearningmastery.com/how-to-investigate-machine-learning-algorithm-behavior/}
\section{引言}
机器学习算法都是一个个复杂的体系，需要通过研究来理解。学习算法的静态描述是一个好的开始，但是这并不足以使我们理解算法的行为，我们需要在动态中来理解算法。

机器学习算法的运行实验，会使你对于不同类型问题得出的实验结论，并对实验结论与算法参数两者的因果关系有一个直观认识。

在这篇文章中，你将会知道怎么研究学习一个机器学习算法。你将会学到5个简单步骤，你可以用来设计和完成你的第一个机器学习算法实验

你会发现机器学习实验不光是学者们的专利，你也可以；你也会知道实验是通往精通的必经之路，因为你可以从经验中学到因果关系的知识， 这是其它地方学不到的。

\section{什么是研究机器学习算法}
当研究一个机器学习算法的时候，你的目标是找到可得到好结果的机器算法行为，这些结果是可以推广到多个问题或者多个类型的问题上。

你通过对算法状态做系统研究来研究学习机器学习算法。这项工作通过设计和运行可控实验来完成

一旦你完成了一项实验，你可以对结论作出解释和提交。这些结论会让你得以管窥在算法变化中因果关系。这就是算法行为和你获得的结论间的关系。

\section{怎样研究学习机器学习算法}
在这一部分，我们将学到5个简单的步骤，你可以通过它来研究学习一个机器算法
\subsection{选择一个算法}
选择一个你有疑问的算法

这个算法可能是你正在某个问题上应用的，或者你发现在其他环境中表现很好，将来你想使用

就实验的意图来说，使用现成的算法是有帮助的。这会给你一个底线：存在bug几率最低

自己实现一个算法可能是了解算法过程的一个好的方式，但是，实验期间，会引入额外的变量，比如bug，和大量必须为算法所做的微观决策
\subsection{确定一个问题}
你必须有一个你试图寻找答案的研究问题。问题越明确，问题越有用

给出的示例问题包括以下几个方面：

KNN算法中，作为样本空间中的一部分的K值在增大时有什么影响？

在SVM算法中，选择不同的核函数在二分类问题上有什么影响 ？

在二分类问题中，逻辑回归上的不同参数的缩放有什么影响 ？

在随机森林模型中，在训练集上增加任意属性对在分类准确性上有什么影响？

针对算法，设计你想回答的问题。仔细考虑，然后列出5个逐渐演变的问题，并且深入推敲那个最精确的
\subsection{设计实验}
从你的问题中挑选出关键元素然后组成你的实验内容。 例如，拿上面的示例问题为例：“二元分类问题中逻辑回归上的不同的参数缩放有什么影响？”

你从这个问题中挑出来用来设计实验的元素是：

属性缩放法：你可以采用像正态化、标准化，将某一属性提升至乘方、取对数等方法

逻辑回归：你想使用哪种已经实现的逻辑回归。

二元分类问题：存在数值属性不同的二分类问题标准。需要准备多种问题，其中一些问题的规模是相同的（像电离层），然而其他一些问题的属性有不同的缩放值（像糖尿病问题）。

性能： 类似分类准确性的模型性能分数是需要的

花时间仔细挑选你问题中的组成元素以便为你的问题给出最佳解答。
\subsection{进行试验并且报告你的结论}
完成你的实验

如果算法是随机的，你需要多次重复实验操作并且记录一个平均数和标准偏差

如果你试图寻找在不同实验（比如带有不同的参数）之间结果的差异，你可能想要使用一种统计工具来标明差异是否统计上显著的（就像学生的t检验）

一些工具像R和scikit-learn/SciPy完成这些类型的实验，但是你需要把它们组合在一起，并且为实验写脚本。其他工具像Weka带有图形用户界面，你所使用的工具不要影响问题和你实验设计的严密

总结你的实验结论。你可能想使用图表。单独呈现结果是不够的，他们只是数字。你必须将数字和问题联系起来，并且通过你的实验设计提取出它们的意义

对实验问题来说，实验结果又暗示着什么呢？

保持怀疑的态度。你的结论上有留什么样的漏洞和局限呢。不要逃避这一部分。知道局限性和知道实验结果一样重要
\subsection{重复}
重复操作

继续研究你选择的算法。你甚至想要重复带有不同参数或者不同的测试数据集的同一个实验。你可能想要处理你试验中的局限性

不要只停留在一个算法上，开始建立知识体系和对算法的直觉

通过使用一些简单工具，提出好的问题，保持严谨和怀疑的态度，你对机器算法行为的理解很快就会到达世界级的水平
\section{研究学习算法不仅仅是学者才能做的}
你也可以学习研究机器学习算法。

你不需要一个很高的学位，你不需要用研究的方式训练，你也不需要成为一名学者

对每个拥有计算机和浓厚兴趣的人来说，机器学习算法的系统研究学习是开放的。事实上，如果你主修机器学习，你一定会适应机器学习算法的系统研究。知识根本不会自己出来，你需要靠自己的经验去得到

当谈论你的发现的适用性时，你需要保持怀疑和谨慎

你不一定提出独一无二的问题。通过研究一般的问题，你也将会收获很多，例如根据一些一般的标准数据集总结出一个参数的普遍影响。你保不住会发现某些具有最优方法的常例的局限性甚至反例。

\section{行动步骤}
在本篇文章中，通过可控实验你知道了研究学习机器学习算法行为的重要性。你掌握了简单的5个步骤，你可以在一个机器学习算法上设计和运行你的第一项实验

采取行动。使用你在这篇博文中学到的步骤，来完成你的第一个机器学习实验。一旦你完成了一个，甚至是很小的一个，你将会获得自信，工具、能力来完成第二个以及更多

我很乐意听到你第一个实验的消息。留下评论，分享你的结论、你的收获

\chapter{框架比较}
\input{chapter\_框架比较.tex}

\chapter{深度学习实例}

\bibliographystyle{plain}
\bibliography{Referenzarchiv}
\end{document} 